% Part: first-order-logic
% Chapter: tableaux
% Section: derivations

\documentclass[../../../include/open-logic-section]{subfiles}

\begin{document}

\iftag{FOL}
      {\olfileid{fol}{tab}{der}}
      {\olfileid{pl}{tab}{der}}

\olsection{真值树}

\begin{explain}
前面已经介绍了什么是假设，也给出了推理规则。分析表就是由这些内容递归生成的：每个分析表要么是仅由一个或多个假设构成的单个分支，要么是通过在某个分支上应用一条推理规则，从一个已有的分析表得到的。
\end{explain}

\begin{defn}[分析表]
关于假设 $\sFmla{S_1}{!A_1}$, \dots,
$\sFmla{S_n}{!A_n}$（其中每个 $S_i$ 是 $\True$ 或~$\False$）的分析表，是一个满足以下条件的带号公式的有限树：
\begin{enumerate}
\item 树最顶部的 $n$ 个带号公式是
  $\sFmla{S_i}{!A_i}$，自上而下排列。
\item 树中每一个不是假设的带号公式，都是由对其上方分支中的带号公式正确应用某条推理规则得到的。
\end{enumerate}
分析表的分支是\emph{封闭的}，当且仅当它同时包含 $\sFmla{\True}{!A}$ 和~$\sFmla{\False}{!A}$，否则是\emph{开放的}。所有分支都封闭的分析表，称为（关于其假设集的）\emph{封闭分析表}。如果一个分析表不封闭，即它含有至少一个开放分支，则它是\emph{开放的}。
\end{defn}

\begin{ex}
  任何假设集自身就是一个分析表，但它通常不是封闭的。（显然，只有当假设中已经包含一对带号公式 $\sFmla{\True}{!A}$ 和~$\sFmla{\False}{!A}$ 时，它才是封闭的。）

  我们可以通过对其中的某个带号公式~$!A$ 应用一条推理规则，从一个（开放或封闭的）分析表得到一个新的、更大的分析表。对带号公式~$!A$ 应用规则时，该规则会将一个或多个新的带号公式添加到包含该~$!A$ 出现的所有分支的末端。

  例如，考虑假设 $\sFmla{\True}{!A \land \lnot
    !A}$。以下就是仅由该假设构成的（开放）分析表：
  \begin{center}
    \begin{tableau}{}
      [\sFmla{\True}{\formula{A} \land \lnot \formula{A}}, just=\TAss]
    \end{tableau}{}
  \end{center}
  我们通过对该假设应用 $\TRule{\True}{\land}$ 规则，得到一个新的分析表。该规则允许我们向分析表添加两行新的带号公式：$\sFmla{\True}{!A}$ 和 $\sFmla{\True}{\lnot !A}$：
  \begin{center}
    \begin{tableau}{}
      [\sFmla{\True}{\formula{A} \land \lnot \formula{A}}, just=\TAss
        [\sFmla{\True}{\formula{A}}, just={\TRule{\True}{\land}[1]},
          [\sFmla{\True}{\lnot \formula{A}}, just={\TRule{\True}{\land}[1]}
          ]
        ]
      ]
    \end{tableau}{}
  \end{center}
  在书写分析表时，我们会在右侧记录所应用的规则（例如，$\TRule{\True}{\land} 1$ 表示该行的带号公式是由对第~$1$ 行的带号公式应用 $\TRule{\True}{\land}$ 规则得到的）。
  现在，这个新分析表包含了额外的带号公式，但只能对其中的一个（即 $\sFmla{\True}{\lnot !A}$）应用规则（本例中是 $\TRule{\True}{\lnot}$ 规则），由此得到如下的封闭分析表：
  \begin{center}
    \begin{tableau}{}
      [\sFmla{\True}{\formula{A} \land \lnot \formula{A}}, just=\TAss
        [\sFmla{\True}{\formula{A}}, just={\TRule{\True}{\land}[1]}
          [\sFmla{\True}{\lnot \formula{A}}, just={\TRule{\True}{\land}[1]}
            [\sFmla{\False}{\formula{A}}, just={\TRule{\True}{\lnot}[3]}, close]
          ]
        ]
      ]
    \end{tableau}{}
  \end{center}
\end{ex}

\end{document}